%\usepackage[small,sf,bf]{titlesec}
\PassOptionsToPackage{table,x11names,dvipsnames,rgb}{xcolor}

\usepackage[utf8]{inputenc}
\usepackage[british]{babel}

\usepackage{xstring}

\usepackage{tikz}
\usetikzlibrary{arrows,shapes,automata,positioning,trees,shadows,mindmap,decorations}
%% From <http://tex.stackexchange.com/questions/137357/how-to-draw-an-arrow-with-two-colors>
%% <http://tex.stackexchange.com/a/137438>
\usetikzlibrary{fadings,decorations.pathmorphing}

\makeatletter
\newif\iftikz@shading@path

\tikzset{
    % There are three circumstances in which the fading sep is needed:
    % 1. Arrows which do not update the bounding box (which is most of them).
    % 2. Line caps/joins and mitres that extend outside the natural bounding 
    %    box of the path (these are not calculated by PGF).
    % 3. Other reasons that haven't been anticipated.
    shading xsep/.store in=\tikz@pathshadingxsep,
    shading ysep/.store in=\tikz@pathshadingysep,
    shading sep/.style={shading xsep=#1, shading ysep=#1},
    shading sep=0.0cm,
}

\def\tikz@shadepath#1{% 
    % \tikz@addmode installs the `modes' (e.g., fill, draw, shade) 
    % to be applied to the path. It isn't usualy for doing more
    % changes to the path's construction.
    \iftikz@shading@path%
    \else%
        \tikz@shading@pathtrue%
        % Get the current path.
        \pgfgetpath\tikz@currentshadingpath%
        % Get the shading sep without setting any other keys.
        \begingroup%
            \pgfsys@beginscope% <- may not be necessary
            \tikzset{#1}%
            \xdef\tikz@tmp{\noexpand\def\noexpand\tikz@pathshadingxsep{\tikz@pathshadingxsep}%
                \noexpand\def\noexpand\tikz@pathshadingysep{\tikz@pathshadingysep}}%
            \pgfsys@endscope%
        \endgroup
        \tikz@tmp%
        % Get the boudning box of the current path size including the shading sep
        \pgfextract@process\pgf@shadingpath@southwest{\pgfpointadd{\pgfqpoint{\pgf@pathminx}{\pgf@pathminy}}%
            {\pgfpoint{-\tikz@pathshadingxsep}{-\tikz@pathshadingysep}}}%%
        \pgfextract@process\pgf@shadingpath@northeast{\pgfpointadd{\pgfqpoint{\pgf@pathmaxx}{\pgf@pathmaxy}}%
            {\pgfpoint{\tikz@pathshadingxsep}{\tikz@pathshadingysep}}}%
        % Clear the path
        \pgfsetpath\pgfutil@empty%                          
        % Save the current drawing mode and options.
        \let\tikz@options@saved=\tikz@options%
        \let\tikz@mode@saved=\tikz@mode%
        \let\tikz@options=\pgfutil@empty%
        \let\tikz@mode=\pgfutil@empty%
        % \tikz@options are processed later on.
        \tikz@addoption{%
            \pgfinterruptpath%
            \pgfinterruptpicture%
                \begin{tikzfadingfrompicture}[name=.]
                \pgfscope%
                    \tikzset{shade path/.style=}% Make absolutely sure shade path is not inherited.
                    \path \pgfextra{%
                        % Set the softpath. Any transformations,draw=none} in #1 will have no effect.
                        % This will *not* update the bounding box...
                        \pgfsetpath\tikz@currentshadingpath%
                        % ...so it is done manually.
                        \pgf@shadingpath@southwest
                        \expandafter\pgf@protocolsizes{\the\pgf@x}{\the\pgf@y}%
                        \pgf@shadingpath@northeast%
                        \expandafter\pgf@protocolsizes{\the\pgf@x}{\the\pgf@y}%
                        % Install the drawing modes and options.
                        \let\tikz@options=\tikz@options@saved%
                        \let\tikz@mode=\tikz@mode@saved%
                    };
                    % Now get the bounding box of the picture.
                    \xdef\pgf@shadingboundingbox@southwest{\noexpand\pgfqpoint{\the\pgf@picminx}{\the\pgf@picminy}}%
                    \xdef\pgf@shadingboundingbox@northeast{\noexpand\pgfqpoint{\the\pgf@picmaxx}{\the\pgf@picmaxy}}%
                    \endpgfscope
                \end{tikzfadingfrompicture}%
            \endpgfinterruptpicture%
            \endpgfinterruptpath%
            % Install a rectangle that covers the shaded/faded path picture.
            \pgftransformreset%
            \pgfpathrectanglecorners{\pgf@shadingboundingbox@southwest}{\pgf@shadingboundingbox@northeast}%
            %
            % Reset all modes.
            \let\tikz@path@picture=\pgfutil@empty%
            \tikz@mode@fillfalse%
            \tikz@mode@drawfalse%
            \tikz@mode@tipsfalse%
            \tikz@mode@doublefalse%
            \tikz@mode@clipfalse%
            \tikz@mode@boundaryfalse%
            \tikz@mode@fade@pathfalse%
            \tikz@mode@fade@scopefalse%
            % Now install shading options.
            \tikzset{#1}%
            \tikz@mode%
            % Make the fading happen.
            \def\tikz@path@fading{.}%
            \tikz@mode@fade@pathtrue%
            \tikz@fade@adjustfalse%
            % Shift the fading to the mid point of the rectangle
            \pgfpointscale{0.5}{\pgfpointadd{\pgf@shadingboundingbox@southwest}{\pgf@shadingboundingbox@northeast}}%
            \edef\tikz@fade@transform{shift={(\the\pgf@x,\the\pgf@y)}}%
            \pgfsetfading{\tikz@path@fading}{\tikz@do@fade@transform}%
            \tikz@mode@fade@pathfalse%              
        }%
    \fi%
}
\tikzset{
    shade path/.code={%
        \tikz@addmode{\tikz@shadepath{#1}}%
    }
}

% requires tikz
\usepackage{colortbl}

\usepackage{float}
\usepackage{graphicx}
\usepackage[normalem]{ulem}

%% use as
%%     \vertcenterimage{\includegraphics{*}}
\newcommand{\vertcenterimage}[1]{\raisebox{-.5\height}{#1}}

%% use as
%%     \flipbox{\includegraphics{*}}
\newcommand{\flipbox}[1]{\scalebox{1}[-1]{#1}}

\usepackage{amsmath}
\usepackage{amsfonts}
\usepackage{bm} % bold mathematics (\bm command)

% set-builder notation using \Set{ ... | ... }
\usepackage{braket}

\usepackage{mathtools}
\usepackage{breqn}

%% make everything in the box bold; both the text and mathematics: using
%% \boldmath from `bm` package
\newcommand{\be}{\bfseries\boldmath}

\usepackage{pifont}% http://ctan.org/pkg/pifont
\newcommand{\cmark}{\ding{51}}%
\newcommand{\xmark}{\ding{55}}%

\newcommand{\todo}[1]{%
	\textcolor{red}{TODO: #1}%
}
\newcommand{\todofig}[1]{%
	\textcolor{red}{TODO figure: \nolinkurl{#1}}%
}


\usetheme{Ilmenau}
\usecolortheme{beaver}
\usepackage{ifxetex}

\ifxetex
	% set font to Tahoma
	\usefonttheme{professionalfonts} % using non standard fonts for beamer
	\usefonttheme{serif} % default family is serif
	\usepackage{fontspec}
	\setmainfont{Tahoma}
\else
	\usepackage[T1]{fontenc}
\fi

\definecolor{CBLRed}{RGB}{205,0,0}

% Title on title slide
\setbeamerfont{title}{size = \Large}
\setbeamercolor{title}{fg = black, bg = white}

\setbeamercolor{frametitle}{fg=white,bg=CBLRed}
\setbeamercolor{section in head/foot}{fg=white,bg=CBLRed}
\setbeamercolor{subsection in head/foot}{fg=white,bg=CBLRed}
\setbeamercolor{author in head/foot}{fg=white,bg=CBLRed}
\setbeamertemplate{headline}{}
\beamertemplatenavigationsymbolsempty

\logo{}
\def\insertlogo
{%
	\color{gray}
	\begin{tabular*}{\paperwidth}{p{0.01\paperwidth}p{0.84\paperwidth}p{0.15\paperwidth}}
		&
		\parbox[c][][c]{\textwidth}{\raggedright
			\vertcenterimage{\includegraphics[width=\textwidth]{gfx/cbl-footer.png}}
		}
		&
		\parbox[c][][c]{\textwidth}{%\raggedleft
			{\large\insertframenumber} %/\inserttotalframenumber
		}
	\end{tabular*}
 }

%\usefonttheme[onlymath]{serif}
\setbeamertemplate{frametitle}[default][center]
\setbeamertemplate{footline}{}
%\setbeamertemplate{footline}[frame number]

%% Change beamer bullets to circles rather than the ball default
\setbeamertemplate{itemize items}[circle]
\setbeamertemplate{enumerate items}[circle]


\usepackage{textcomp}
\usepackage{fancyvrb}
\usepackage{changepage}
\usepackage{multicol}
\usepackage{wasysym}
\usepackage{listings}

\lstset{%basicstyle=\small\ttfamily,
%numbers=left,
%escapeinside=||
}
\newenvironment{indented}{\begin{adjustwidth}{1.5em}{}}{\end{adjustwidth}}

% http://tex.stackexchange.com/questions/12550/changing-default-width-of-blocks-in-beamer/12551#12551
\newenvironment<>{varblock}[2][.9\textwidth]{%
  \setlength{\textwidth}{#1}
  \begin{actionenv}#3%
    \def\insertblocktitle{#2}%
    \par%
    \usebeamertemplate{block begin}}
  {\par%
    \usebeamertemplate{block end}%
  \end{actionenv}}


%% TOC at beginning of sections/subsections
%\AtBeginSection[]
%{
%% <beamer>{Table of Contents}
%        \begin{frame}
%                \tableofcontents[
%                                currentsection,
%                                hideothersubsections,
%                                sectionstyle=show/shaded,
%                                subsectionstyle=hide/hide,
%                ]
%        \end{frame}
%}
%\AtBeginSubsection[]
%{
%   \begin{frame}
%                \begin{tabular}{p{0.3\paperwidth}p{0.7\paperwidth}}
%                        \tableofcontents[
%                                        currentsection,
%                                        hideothersubsections,
%                                        sectionstyle=show/shaded,
%                                        subsectionstyle=hide/hide,
%                        ]
%                &
%                        \begin{varblock}[0.5\paperwidth]{}
%                                \vspace{.2cm}
%                                \tableofcontents[
%                                        currentsubsection,
%                                        hideothersubsections,
%                                        sectionstyle=hide/hide,
%                                        subsectionstyle=show/shaded,
%                                ]
%                        \end{varblock}
%                \end{tabular}
%   \end{frame}
%}


\AtBeginSubsection[]
{
	\begin{frame}
		\begin{multicols}{2}
			\tableofcontents[
				currentsection,
				hideothersubsections,
				sectionstyle=show/shaded,
				subsectionstyle=show/shaded,
			]
		\end{multicols}
  \end{frame}
}

%% TikZ arrows
%% From <https://tex.stackexchange.com/questions/61507/drawing-arrows-in-beamer>
\tikzset{
    myarrow/.style={
        draw,
        fill=orange,
        single arrow,
        minimum height=3.5ex,
        single arrow head extend=1ex
    }
}
\newcommand{\arrowup}{%
\tikz [baseline=-0.5ex]{\node [myarrow,rotate=90] {};}
}
\newcommand{\arrowdown}{%
\tikz [baseline=-1ex]{\node [myarrow,rotate=-90] {};}
}
\newcommand{\arrowright}{%
\tikz [baseline=-0.5ex]{\node [myarrow,rotate=0] {};}
}
\newcommand{\arrowleft}{%
\tikz [baseline=-0.5ex]{\node [myarrow,rotate=180] {};}
}
